% GENERATED FILE. Edit canonical lessons, then run npm run generate:books.
% canonical-source-hash: fnv1a64:57718578d5471c86
% canonical-lessons: PT-C508-arbitro, PT-C508-apito, PT-C508-apitar, PT-C508-bancada, PT-C508-suplente

\chapter{O árbitro e o apito}
\label{ch:pt-o-arbitro-e-o-apito}
\hlchaptermodality{508}

\begin{chapteropening}
\textbf{By the end of this chapter:} \emph{I can talk about the referee, the stands and the bench.}

Complete the last lesson of chapter 508: I can talk about the referee, the stands and the bench.
\end{chapteropening}

\section[o árbitro]{o árbitro — referee}

\label{lesson:PT-C508-arbitro}



\begin{quote}
Before the new one: say the Portuguese for (long) walk, then the Portuguese for mountaineering.
\end{quote}

\subsection*{You'll want to know: o árbitro}
\textbf{o árbitro} — \textquotedblleft{}referee\textquotedblright{}.

\begin{quote}
O árbitro mostrou um cartão amarelo.
\end{quote}

\begin{quote}
\emph{The referee showed a yellow card.}
\end{quote}

A woman is \emph{a árbitra}.

\begin{cousinweb}
\emph{árbitro} is a learned borrowing from Latin \emph{arbiter} \textquotedblleft{}a judge, an umpire\textquotedblright{} (EtymDB).
\end{cousinweb}

\begin{grammarlens}[title={What you've built}]
One more word for telling, explaining and describing at B1.
\end{grammarlens}

\subsection*{Your turn}
\begin{itemize}
\raggedright
  \item \emph{Say it:} \emph{o árbitro}
  \item \emph{Say it:} \emph{o árbitro}, once more
  \item \emph{Say it:} a sentence of your own with \emph{o árbitro}
  \item \emph{Recall:} say the Portuguese for (long) walk, then the Portuguese for mountaineering, then say \emph{o árbitro} again
\end{itemize}

\begin{tcolorbox}[breakable,colback=teal!4,colframe=teal!35!black,title={Before you move on}]
What does o árbitro mean? (\textquotedblleft{}Referee\textquotedblright{}.) Say it once more.
\end{tcolorbox}
\section[o apito]{o apito — whistle}

\label{lesson:PT-C508-apito}



\begin{quote}
Before the new one: say the Portuguese for mountaineering, then the Portuguese for referee.
\end{quote}

\subsection*{You'll want to know: o apito}
\textbf{o apito} — \textquotedblleft{}whistle\textquotedblright{}.

\begin{quote}
Ouviu-se o apito final.
\end{quote}

\begin{quote}
\emph{The final whistle was heard.}
\end{quote}

The sound itself and the object.

\begin{cousinweb}
\emph{apito} is probably imitative of a sharp whistling sound; the origin is not settled.
\end{cousinweb}

\begin{grammarlens}[title={What you've built}]
One more word for telling, explaining and describing at B1.
\end{grammarlens}

\subsection*{Your turn}
\begin{itemize}
\raggedright
  \item \emph{Say it:} \emph{o apito}
  \item \emph{Say it:} \emph{o apito}, once more
  \item \emph{Say it:} a sentence of your own with \emph{o apito}
  \item \emph{Recall:} say the Portuguese for mountaineering, then the Portuguese for referee, then say \emph{o apito} again
\end{itemize}

\begin{tcolorbox}[breakable,colback=teal!4,colframe=teal!35!black,title={Before you move on}]
What does o apito mean? (\textquotedblleft{}Whistle\textquotedblright{}.) Say it once more.
\end{tcolorbox}
\section[apitar]{apitar — to blow the whistle; to referee}

\label{lesson:PT-C508-apitar}



\begin{quote}
Before the new one: say the Portuguese for referee, then the Portuguese for whistle.
\end{quote}

\subsection*{You'll want to know: apitar}
\textbf{apitar} — \textquotedblleft{}to blow the whistle; to referee\textquotedblright{}.

\begin{quote}
O árbitro apitou o fim do jogo.
\end{quote}

\begin{quote}
\emph{The referee blew the whistle for the end of the match.}
\end{quote}

Of a machine: to beep. \emph{O micro-ondas apitou}.

\begin{cousinweb}
\emph{apitar} is built on \emph{apito} \textquotedblleft{}whistle\textquotedblright{}, which is probably imitative; the origin is not settled.
\end{cousinweb}

\begin{grammarlens}[title={What you've built}]
One more word for telling, explaining and describing at B1.
\end{grammarlens}

\subsection*{Your turn}
\begin{itemize}
\raggedright
  \item \emph{Say it:} \emph{apitar}
  \item \emph{Say it:} \emph{apitar}, once more
  \item \emph{Say it:} a sentence of your own with \emph{apitar}
  \item \emph{Recall:} say the Portuguese for referee, then the Portuguese for whistle, then say \emph{apitar} again
\end{itemize}

\begin{tcolorbox}[breakable,colback=teal!4,colframe=teal!35!black,title={Before you move on}]
What does apitar mean? (\textquotedblleft{}To blow the whistle; to referee\textquotedblright{}.) Say it once more.
\end{tcolorbox}
\section[a bancada]{a bancada — stand (in a stadium)}

\label{lesson:PT-C508-bancada}



\begin{quote}
Before the new one: say the Portuguese for whistle, then the Portuguese for to blow the whistle.
\end{quote}

\subsection*{You'll want to know: a bancada}
\textbf{a bancada} — \textquotedblleft{}stand (in a stadium)\textquotedblright{}.

\begin{quote}
Os adeptos cantavam na bancada.
\end{quote}

\begin{quote}
\emph{The fans were singing in the stand.}
\end{quote}

Also a kitchen worktop: \emph{a bancada da cozinha}.

\begin{cousinweb}
\emph{bancada} is built on \emph{banco} \textquotedblleft{}bench\textquotedblright{}, which goes back to a Germanic word for a bench (EtymDB).
\end{cousinweb}

\begin{grammarlens}[title={What you've built}]
One more word for telling, explaining and describing at B1.
\end{grammarlens}

\subsection*{Your turn}
\begin{itemize}
\raggedright
  \item \emph{Say it:} \emph{a bancada}
  \item \emph{Say it:} \emph{a bancada}, once more
  \item \emph{Say it:} a sentence of your own with \emph{a bancada}
  \item \emph{Recall:} say the Portuguese for whistle, then the Portuguese for to blow the whistle, then say \emph{a bancada} again
\end{itemize}

\begin{tcolorbox}[breakable,colback=teal!4,colframe=teal!35!black,title={Before you move on}]
What does a bancada mean? (\textquotedblleft{}Stand (in a stadium)\textquotedblright{}.) Say it once more.
\end{tcolorbox}
\section[o suplente]{o suplente — substitute (player)}

\label{lesson:PT-C508-suplente}



\begin{quote}
Before the new one: say the Portuguese for to blow the whistle, then the Portuguese for stand.
\end{quote}

\subsection*{You'll want to know: o suplente}
\textbf{o suplente} — \textquotedblleft{}substitute (player)\textquotedblright{}.

\begin{quote}
O suplente entrou e marcou um golo.
\end{quote}

\begin{quote}
\emph{The substitute came on and scored a goal.}
\end{quote}

\emph{O banco de suplentes} — the bench.

\begin{cousinweb}
\emph{suplente} is a learned borrowing from Latin \emph{supplēns}, from \emph{supplēre} \textquotedblleft{}to fill up, to make good\textquotedblright{} (Lewis \& Short).
\end{cousinweb}

\begin{grammarlens}[title={What you've built}]
One more word for telling, explaining and describing at B1.
\end{grammarlens}

\subsection*{Your turn}
\begin{itemize}
\raggedright
  \item \emph{Say it:} \emph{o suplente}
  \item \emph{Say it:} \emph{o suplente}, once more
  \item \emph{Say it:} a sentence of your own with \emph{o suplente}
  \item \emph{Recall:} say the Portuguese for to blow the whistle, then the Portuguese for stand, then say \emph{o suplente} again
\end{itemize}

\begin{tcolorbox}[breakable,colback=teal!4,colframe=teal!35!black,title={Before you move on}]
What does o suplente mean? (\textquotedblleft{}Substitute (player)\textquotedblright{}.) Say it once more.
\end{tcolorbox}
